\documentclass[11pt,a4paper]{article}
\usepackage{../shared/luxcover}
\usepackage[utf8]{inputenc}
\usepackage[T1]{fontenc}
\usepackage{amsmath,amssymb,amsfonts,amsthm}
\usepackage{hyperref}
\usepackage{booktabs}
\usepackage{enumitem}
\usepackage{xcolor}
\usepackage{graphicx}

\hypersetup{colorlinks=true,linkcolor=black,citecolor=blue,urlcolor=blue}

\newtheorem{definition}{Definition}[section]
\newtheorem{proposition}{Proposition}[section]
\newtheorem{remark}{Remark}[section]

\newcommand{\bps}{\,\text{bps}}
\newcommand{\dev}{\textsuperscript{\textsf{\scriptsize dev}}}
\newcommand{\sither}{\textsuperscript{\textsf{\scriptsize sim}}}

\title{Unified V2/V3/V4 Price Coherence:\\
A CEX-Oracle Market-Maker with Arbitrage-Driven\\
Convergence and LP-Fee Recycling on a Single Pool Precompile}

\author{
Zach Kelling\thanks{Corresponding author: \texttt{zach@lux.network}}\\
\textit{Lux Industries, Inc.}\\
\texttt{research@lux.network}
}

\date{}

\begin{document}

\luxcoverpage

\begin{abstract}
On most decentralized exchanges the constant-product (V2), concentrated-liquidity
(V3) and order-book/V4 venues are \emph{separate contracts} whose prices are kept
loosely aligned only by external arbitrage. We describe and measure a different
design realized on the Lux DEX: V2, V3 and V4 are three \emph{liquidity modes of a
single pool} hosted by one EVM precompile (\texttt{0x9999}), sharing one pool
identity, one fee tier and one maker peg. A single market-maker tracks a live
centralized-exchange (CEX) mid (Kraken public market data) and pegs the book to it;
an arbitrage loop harvests the residual requote-lag gap, pulling the on-chain price
back to the reference. Because the LP is also the arber on a sovereign L1, only the
EIP-1559 base fee is a true cost, so the loop recycles fee income to the dominant LP
net of a small burn. We implement the CEX oracle and the coherence/arbitrage harness,
and drive real fills on the Lux devnet \texttt{0x9999} \texttt{PoolManager} (chain
\texttt{96367}). We report exact on-chain costs---\texttt{swap} $73{,}608$ gas,
a maker requote (\texttt{deposit}+\texttt{place}) $114{,}272$ gas, $25$~gwei base
fee---and establish the CLOB fee mechanism empirically (a parity fill of size $10$
returns exactly $10$: the V4 CLOB charges \emph{no} taker fee; the LP earns the quoted
spread, while the $30\bps$ tier applies to the V2/V3 curve). A simulation calibrated by
these measurements and driven by a real $12$-hour Kraken ETH/USD trace
($7.66\bps$ per-minute volatility) shows the three modes track the oracle within
$0.3$--$13\bps$ for requote intervals of $2$\,s--$10$\,min, with cross-version
dispersion below $14\bps$---inside the $30\bps$ no-arbitrage band. The cycle is
self-sustaining at every interval; the burn is minimized ($\sim\!\$100$--$160$/day for
three pools) at a $60$--$300$\,s requote band where requote-gas and arbitrage-gas
balance, with a break-even external volume of $\sim\!\$33$--$53$k/day. We give the
no-arbitrage-band model relating requote interval, volatility, fee and gas to tracking
error and break-even volume, and discuss oracle-manipulation, MEV and inventory risks.
\end{abstract}

\section{Introduction}

A decentralized exchange that lists an asset also traded on centralized venues faces a
standing problem: its on-chain price must stay aligned with the broader market, or it
becomes a faucet for arbitrageurs at the liquidity providers' expense. The
conventional remedy is passive---deploy automated market-maker (AMM) pools and let
external arbitrage drag the price toward the global mid. This is robust but leaks value
to outside arbitrageurs and provides no guarantee of \emph{coherence} across the
several pool types (V2 constant-product, V3 concentrated liquidity, V4/order-book) a
mature DEX runs in parallel, each a separate contract with its own price.

This paper studies an \emph{active} design realized on the Lux DEX, where the three
``versions'' are not separate contracts but \emph{three liquidity modes of one pool}
inside a single EVM precompile. A single market-maker, which we operate as the dominant
LP, pegs the shared book to a live CEX reference; an arbitrage loop, which we also
operate, harvests the residual gap that opens during requote lag and pulls the price
back. Because on a sovereign L1 the operator is also the validator, the priority tip
recycles and only the base fee is burned, so the loop recycles fee income to the LP net
of a small, measurable burn.

\paragraph{Contributions.}
\begin{enumerate}[leftmargin=1.4em]
\item \textbf{A unified-pool coherence mechanism} (\S\ref{sec:dex},\S\ref{sec:mech}):
  pegging one maker quote to a CEX oracle keeps V2/V3/V4 coherent \emph{by
  construction}, because they share one pool identity and one peg---not three
  reconciled AMMs.
\item \textbf{An implementation} (\S\ref{sec:impl}): a Kraken public-data oracle as a
  drop-in maker price source, and a coherence/arbitrage harness driving the
  \texttt{0x9999} money path; both build and ship with tests, and the oracle is
  live-proven against the real Kraken API.
\item \textbf{Real on-chain measurements} (\S\ref{sec:results}) on Lux devnet: exact
  per-operation gas, block cadence, latency, and an empirical determination of the
  CLOB fee mechanism (no taker fee; LP earns the spread).
\item \textbf{A calibrated study} (\S\ref{sec:results},\S\ref{sec:econ}): a simulation
  pinned to the measured gas/fee/block parameters and driven by a real Kraken trace,
  quantifying tracking error, cross-version dispersion, and the fee/burn economics, and
  the no-arbitrage-band model that explains them (\S\ref{sec:mech}).
\end{enumerate}

We are careful throughout to label every quantity \dev~(real, from an accepted devnet
receipt or log) or \sither~(from the calibrated model). No figure is a vendor number.

\section{The Lux unified DEX}\label{sec:dex}

The Lux DEX money path is a single EVM precompile at \texttt{0x9999}, the V4-style
\texttt{PoolManager}. A pool is identified by a \texttt{PoolKey} (currency pair, fee
tier, tick spacing, hooks). The precompile hosts three liquidity modes against the same
pool identity:

\begin{center}
\begin{tabular}{@{}llll@{}}
\toprule
``Version'' & Mode on \texttt{0x9999} & Mechanism & Devnet status \\
\midrule
V4 & native CLOB & \texttt{swapPlace}/\texttt{swapDeposit}/\texttt{swap} & live, fills proven \\
V3 & concentrated liquidity & \texttt{modifyLiquidity} (sqrtPrice/ticks) & precompile present \\
V2 & constant product $xy{=}k$ & \texttt{BindAMMPool}$\to$router blend & present; no bind selector \\
\bottomrule
\end{tabular}
\end{center}

Read-only views are separate precompiles, all live on devnet: \texttt{0x9998}
(Quoter, advisory swap quote), \texttt{0x9997} (StateView, pool/book/position read),
and \texttt{0x9012} (LXRouter, which blends the V4 native pool, the V2 constant-product
reserves, and external venues). The router's V2 quote reads a reserve row bound in
\texttt{0x9999} storage; its ``external V3'' path is a stub, because the \emph{native}
V3 equivalent is \texttt{0x9999} concentrated liquidity, not a foreign Uniswap-V3
deployment.

\begin{remark}[Coherence is structural]
Because the three modes share one \texttt{PoolKey} identity and one fee tier, a market
maker that pegs \emph{one} quote pegs all three. There is no cross-contract
reconciliation: the ``unified V2/V3/V4 coherence'' reduces to a single price the maker
controls. This is the central architectural advantage over a three-contract Uniswap
deployment.
\end{remark}

\section{Mechanism}\label{sec:mech}

\paragraph{Setup.} Let $P(t)$ be the live CEX reference mid (Kraken). The maker holds a
peg $P_{\text{peg}}$, refreshed to $P(t)$ every requote interval $T$. The pool's
marginal price $m(t)$ starts at the peg and is moved by trades. We are the dominant LP
and we run the arbitrage loop; external flow is exogenous.

\paragraph{No-arbitrage band.} An arbitrageur closing a price gap $\delta = (m-P)/P$
trades size $s$ (base units) at price $P$. The round trip earns $\delta\, s P$, pays the
fee $f\, s P$ (the $30\bps$ pool tier on V2/V3, or the maker's quoted spread on V4), and
burns gas worth $g = \gamma\,\beta_{\text{base}}\,p_{\textsc{lux}}$ in quote units,
where $\gamma$ is the swap gas, $\beta_{\text{base}}$ the base fee, and $p_{\textsc{lux}}$
the LUX price. It profits iff
\begin{equation}\label{eq:band}
|\delta| \;>\; \beta \;\equiv\; f \;+\; \frac{g}{sP}.
\end{equation}
$\beta$ is the half-width of the no-arbitrage band: the pool price is pinned to $P$
within $\pm\beta$ whenever the arb is active. On a sovereign L1 the tip recycles to the
operator's validator, so $g$ is the base-fee burn only and $g/(sP)\!\ll\!f$; hence
$\beta\!\approx\!f$ (fee-dominated).

\paragraph{Tracking error and the requote/arb trade-off.} Between requotes the reference
drifts while the maker quote is stale; the taker-facing tracking error grows as
$\mathbb{E}|m-P|/P \approx \sigma\sqrt{T}$ for per-unit-time volatility $\sigma$, until a
requote ($P_{\text{peg}}\!\leftarrow\!P$, gas $\gamma_{rq}$ per pool) or an arb (when the
drift exceeds $\beta$) resets it. Thus two controls hold coherence: \emph{fast requoting}
(tight tracking, high requote-gas) and \emph{arbitrage} (caps drift at $\beta$, one swap
per correction). The optimum minimizes total burn for a tolerated tracking error.

\paragraph{Cycle economics.} With the LP and the arber both being us, arb fees and arb
adverse selection are internal washes. Real income is the fee on \emph{external} volume
$V_{\text{ext}}$; real cost is the base-fee burn of requoting and of our arb swaps:
\begin{equation}\label{eq:econ}
\text{net} \;=\; f\,V_{\text{ext}} \;-\; p_{\textsc{lux}}\,\beta_{\text{base}}
\Big( N_{rq}\,\gamma_{rq} + N_{arb}\,\gamma_{arb}\Big),
\end{equation}
with $N_{rq}=\#\text{pools}\cdot \lfloor \mathcal{H}/T\rfloor$ requotes over horizon
$\mathcal{H}$ and $N_{arb}$ arb swaps. The cycle is self-sustaining ($\text{net}>0$)
above the break-even external volume
\begin{equation}\label{eq:breakeven}
V^\star \;=\; \frac{p_{\textsc{lux}}\,\beta_{\text{base}}\,(N_{rq}\gamma_{rq}+N_{arb}\gamma_{arb})}{f}.
\end{equation}

\section{Implementation}\label{sec:impl}

\paragraph{Kraken oracle.} We add a price source reading Kraken's public, unauthenticated
market-data endpoints (\texttt{/0/public/Ticker} for BBO, \texttt{/0/public/Depth} for
book). In the maker it is an aggregator plug-in (the same shape as the existing on-chain
oracle sources) exposed as the live mid the seeder tracks; a composition primitive lets
the CEX feed (listed majors) and the native-AMM feed (LUX/\(*\)) coexist behind the
single feed the seeder consumes. In the arbitrage client it is a scalar-mid reader. The
oracle is hermetically tested and live-proven (a real \texttt{BTC/USD} mid of
\$59{,}507.75 pulled from the public API during testing).

\paragraph{Coherence harness.} The harness drives the full loop on \texttt{0x9999} using
only the proven primitives (\texttt{Initialize}, \texttt{Deposit}, \texttt{Place},
\texttt{Swap}, and the \texttt{DEXFill} log reader). Each round pegs the maker bid to the
oracle mid, has the taker trade it, and records a per-round sample---realized fill price,
implied fee, gas, block, latency---so the economics are \emph{measured}. A drift-injection
window mispriced relative to the reference exercises the arbitrage leg.

\section{Methodology and results}\label{sec:results}

\paragraph{On-chain [dev].} We ran the harness against Lux devnet (chain \texttt{96367},
\texttt{0x9999}, the genesis-funded index-0 account). Gas is exact from receipts; block
timing from block timestamps; fills from \texttt{DEXFill} logs.

\begin{center}
\begin{tabular}{@{}lrr@{}}
\toprule
operation & gas (\dev) & base-fee burn @25\,gwei \\
\midrule
\texttt{swapDeposit} & $41{,}584$ & $0.001040$ LUX \\
\texttt{swapPlace}   & $72{,}688$ & $0.001817$ LUX \\
\texttt{swap} (taker/arb) & $73{,}608$ & $0.001840$ LUX \\
maker requote (\texttt{deposit}+\texttt{place}) & $114{,}272$ & $0.002857$ LUX \\
\bottomrule
\end{tabular}
\end{center}

The devnet builds one transaction per block; warm inter-block time is $2$--$5$\,s, with
$25$--$50$\,s gaps after the builder idles. Base fee is pinned at the $25$~gwei floor.
Taker-swap latency was $\approx\!5$\,s in steady state with cold-start spikes to
$40$--$67$\,s. Six \texttt{DEXFill}s settled to accepted state.

\begin{proposition}[CLOB fee mechanism, empirical]\label{prop:fee}
The V4 CLOB charges no taker fee; the LP's CLOB income is the quoted spread. A taker
selling $10$ base units at parity into a parity bid received exactly $10$ (block 14),
i.e. zero fee. The $30\bps$ \texttt{Fee} tier applies to the V2/V3 curve, not to CLOB
matches.
\end{proposition}

\paragraph{Simulation [sim].} The devnet block-production cadence makes long horizons
impractical to run live, so we drive a calibrated model: V2 ($xy{=}k$), V3 (concentrated,
$20\times$ depth), V4 (CLOB, $5\times$ depth, $1\bps$ tick) tracking the \emph{real}
Kraken ETH/USD $1$-minute trace ($721$ points, $12$\,h, $\sigma_{1\text{m}}\!=\!7.66\bps$,
$221\bps$ realized range). Calibration is the measured $\gamma_{arb}\!=\!73{,}608$,
$\gamma_{rq}\!=\!114{,}272$, $\beta_{\text{base}}\!=\!25$~gwei, block $=2$\,s, with
$f\!=\!30\bps$ and $p_{\textsc{lux}}\!=\!\$12.5$; pool depth, concentration and organic-flow
rate are explicit modelling assumptions. Table~\ref{tab:sweep} sweeps the requote interval.

\begin{table}[h]
\centering
\caption{Requote-interval sweep [sim], $f{=}30\bps$, calibrated by \S\ref{sec:results}.
RMS tracking error per mode, p99 cross-version dispersion, arb count, external fee income
and base-fee burn over the $12$\,h horizon (three pools), and net.}
\label{tab:sweep}
\begin{tabular}{@{}rrrrrrrrr@{}}
\toprule
$T$ & V2 RMS & V3 RMS & V4 RMS & disp.\ p99 & arbs & fee & burn & net \\
(s) & (bps) & (bps) & (bps) & (bps) & & (\$) & (LUX) & (\$) \\
\midrule
2   & 0.69 & 0.26 & 0.29 & 1.90 & 0 & 15{,}370 & 185.1 & 13{,}056 \\
6   & 1.06 & 0.55 & 0.58 & 1.90 & 0 & 15{,}370 & 61.7 & 14{,}598 \\
12  & 1.56 & 1.00 & 1.03 & 3.80 & 0 & 15{,}370 & 30.9 & 14{,}984 \\
30  & 2.91 & 2.33 & 2.35 & 3.80 & 0 & 15{,}370 & 12.3 & 15{,}216 \\
60  & 5.05 & 4.49 & 4.52 & 5.70 & 88 & 15{,}370 & 6.3 & 15{,}291 \\
120 & 7.27 & 6.71 & 6.73 & 7.60 & 522 & 15{,}370 & 4.1 & 15{,}319 \\
300 & 11.21 & 10.42 & 10.46 & 11.40 & 1{,}453 & 15{,}370 & 3.9 & \textbf{15{,}321} \\
600 & 13.40 & 12.82 & 12.83 & 13.31 & 2{,}633 & 15{,}370 & 5.5 & 15{,}302 \\
\bottomrule
\end{tabular}
\end{table}

Four observations. (i)~Tracking error grows monotonically with $T$, from $\sim\!0.3\bps$
at $2$\,s to $\sim\!13\bps$ at $10$\,min, as Eq.~\eqref{eq:band}--model predicts.
(ii)~V3 (concentrated) tracks tightest and V2 (shallow) loosest, the spread being small
($\sim\!0.3$--$0.6\bps$): deeper liquidity resists organic-flow impact. (iii)~Cross-version
dispersion stays well inside the $30\bps$ band ($1.9\!\to\!13.3\bps$ p99): coherence holds
because one maker pegs all three. (iv)~Arbitrage fires only once $T$ is large enough that
drift exceeds $\beta$ ($\geq\!60$\,s here), at which point it---not requoting---is the
corrective force.

\section{Economics: is the cycle self-sustaining?}\label{sec:econ}

By Eq.~\eqref{eq:econ}, income is the external fee and cost is the base-fee burn of
requoting plus arbitrage. The two burn terms trade off: requote burn dominates at small
$T$ ($185$ LUX at $2$\,s), arb burn at large $T$; total burn is minimized in a
$60$--$300$\,s band at $\sim\!4$--$6$ LUX over $12$\,h ($\sim\!\$100$--$160$/day for three
pools). Against the modelled external fee income ($\sim\!\$30$k/day) this is a
$200$--$300\times$ margin, and net is positive at \emph{every} interval. By
Eq.~\eqref{eq:breakeven} the break-even external volume falls from $\$257$k/day (2\,s) to
$\sim\!\$33$k/day in the optimal band.

\begin{proposition}[Self-sustaining]\label{prop:sustain}
Under the calibrated model the cycle is self-sustaining for any external daily volume
above $\sim\!\$50$k; in the burn-optimal $60$--$300$\,s requote band the break-even is
$\sim\!\$33$--$53$k/day.
\end{proposition}

\noindent The policy implication is to \emph{requote at $\sim\!60$--$120$\,s and let the
arbitrage loop hold the no-arbitrage band}, rather than pay for fast requoting: an arb
correction is one $73{,}608$-gas swap, cheaper than re-pegging three pools every few
seconds.

\section{Related work}

Constant-product and concentrated-liquidity AMMs~\cite{uniswapv2,uniswapv3} and the
singleton V4 \texttt{PoolManager}~\cite{uniswapv4} define the venue surface; our pools
are V4-style but host the order book, concentrated liquidity and a constant-product
reserve as modes of one identity. Loss-versus-rebalancing~\cite{lvr} formalizes the
adverse selection an oracle-lagging LP suffers; our no-arbitrage band
(Eq.~\eqref{eq:band}) and self-arb policy are the operational response, and our requote
sweep quantifies the LVR/gas trade-off directly. Oracle-pegged or ``oracle AMM''
designs~\cite{mev} reduce arbitrage leakage by pricing off an external feed; we go
further by making the operator both the dominant LP and the arber on a sovereign L1, so
the leakage is recaptured net of base-fee burn.

\section{Limitations and future work}

The on-chain horizon was short: the devnet builder stalls between transactions (a known
node-version regression), so the long-horizon tracking and economics are from the
calibrated model, not a multi-hour live run; the real measurements pin the model's cost
parameters and its fee mechanism (Prop.~\ref{prop:fee}). The native V2 mode has no
external bind selector today---\texttt{BindAMMPool} is only invoked internally after a
fill---so a permissioned operator-only \texttt{bindAMMPool} entry point is needed to seed
real V2 reserves; V3/V4 are ready on \texttt{0x9999}. The model's pool depth,
concentration factor and organic-flow rate are explicit assumptions; tightening them with
live order-flow data is future work, as is multi-CEX median pricing with staleness and
deviation circuit-breakers (\S below).

\paragraph{Risks.} The peg trusts the CEX feed: a spoofed or halted feed mis-pegs the
book (mitigation: median across feeds, staleness/deviation breakers, per-requote move
clamps). Requote-lag quotes are picked off by external arbitrageurs---this \emph{is} the
tracking error---trading off against requote gas. A trending market fills one side and
accumulates inventory, handled by inventory-skew quoting and size limits. A single
dominant maker pegged to one CEX is a centralization point appropriate for bootstrapping
a market, not an end state.

\section{Conclusion}

Hosting V2, V3 and V4 as liquidity modes of one pool precompile turns ``cross-version
price coherence'' into a single maker peg, which a CEX oracle drives and an arbitrage
loop defends within a fee-width no-arbitrage band. On Lux devnet we measured the exact
costs of the loop and established that the V4 CLOB earns the spread rather than a fee; a
calibrated model driven by a real Kraken trace shows the three modes tracking the oracle
within $0.3$--$13\bps$ and staying coherent within $14\bps$, with the fee-recycling cycle
self-sustaining above a modest external volume. The remaining gaps---a block-production
stall and a missing V2 bind selector---are operational, not architectural.

\begin{thebibliography}{9}
\bibitem{uniswapv2} H.~Adams \emph{et al.} Uniswap v2 Core. 2020.
\bibitem{uniswapv3} H.~Adams \emph{et al.} Uniswap v3 Core. 2021.
\bibitem{uniswapv4} Uniswap Labs. Uniswap v4 Core (singleton \texttt{PoolManager}). 2023.
\bibitem{lvr} J.~Milionis, C.~C.~Moallemi, T.~Roughgarden, A.~L.~Zhang.
  Automated Market Making and Loss-Versus-Rebalancing. 2022.
\bibitem{mev} P.~Daian \emph{et al.} Flash Boys 2.0: Frontrunning, Transaction
  Reordering, and Consensus Instability in Decentralized Exchanges. IEEE S\&P, 2020.
\end{thebibliography}

\end{document}
